\begin{textsample}{NEG Dim 1 – global\_north\_2024\_12 – Score 1.00 – global\_north\_2024\_12/118989544.txt}  \label{ex:f1_neg_222}
A private schools group has apologised afterhosting speakers who allegedly accused Israel of"ethnic cleansing"at an equality conference branded a festival of"Jew hate".
According to an open letter, speakers at last week ' s event in Denver described the establishment of the state of Israel as a"racist"endeavour, accused Israel of genocide, and downplayed Hamas ' s Oct 7 attack in which more than 1,200 people were killed.
Called the People of Color Conference, the annual gathering hosted by the National Association of Independent Schools invited around 1,700 private schools from across the United States, including 60 Jewish day schools.
Some of the country ' s most prestigious private schools sent delegations to the conference, including Manhattan ' s Dalton School, Brearley, Collegiate, Trinity, Fieldston, as well as Los Angeles ' Milken Community Schools. ' Jewish students hid their Star of David necklaces ' The leaders of several prominent Jewish groups have since released an open letter denouncing comments made by speakers at the diversity and inclusion conference, calling their remarks"extreme, biased anti-Zionist and anti-Israel rhetoric"and accusing them of creating a"hostile"environment for Jewish students.
Some Jewish students felt"so unsafe"that they left the conference early and hid their Star of David necklaces as audience members"glared and whispered,"according to the letter."Jewish students and faculty attending PoCC were forced to hear this damaging and anti-Semitic rhetoric repeated time and again and watch as their peers applauded,"according to the letter, co-authored by the Anti-Defamation League, American Jewish Committee, Jewish Federations of North America and Prizmah."These occurrences, along with others reported by Jewish attendees, display a fundamental undermining of the principles of inclusivity and equity that NAIS stands for,"the letter claims."No student should ever be made to feel this way because of their identity."' Jew hate ' Sarah Shulkind, the head teacher of Milken Community School, which was at the conference, described the comments as"Jew hate"."Most horrifying is that the thousands of educators that were in that room teach at the nation ' s most elite schools... they stood and cheered,"she told the New York Post."It reveals to me two things : virulent, unapologetic Jew hate, and profound ignorance."In response to the letter, Debra Wilson, the association ' s president, expressed remorse over the"divisive and hurtful rhetoric"and said that future speakers ' addresses would be vetted."No last-minute changes will be permitted without explicit review and approval,"she said,"That any student would feel the need to conceal their identity at our conference is antithetical to our mission and our values."We are committed to ensuring that future NAIS events will be places where difficult topics and conversations are approached with sensitivity and care for all."Concerns about the Denver conference focused on the remarks of keynote speaker, Dr Suzanne Barakat, an Arab American physician who has been a prominent voice on the subject of anti-Muslim violence.
One of the conference speakers was Dr Suzanne Barakat In her address, Dr Barakat allegedly characterised Zionism as when"some European Jews decided that the solution to solving anti-Semitism in Europe and Russia was the establishment of a state in Palestine"and also accused Israel of"genocide"to applause from the crowd.
According to transcripts obtained by The New York Post, Dr Barakat called the establishment of the Jewish state"ethnic cleansing".
Ruha Benjamin, a professor of African American studies at Princeton University who was referred to in the letter, criticised NAIS ' s decision to apologise and defended calling Israel ' s war in Gaza genocide."The weaponisation of the charge of anti-Semitism is a disservice to everyone,"she wrote in an email seen by The New York Times.
She added that such accusations are"watering down its meaning and wielding it against anyone who dares name the reality that Palestinians are living".
The row comes amid a spate of anti-Semitic incidents reported on college and school campuses over the last year following escalating tensions between administrators and students over the Israel-Hamas war.
A diversity administrator at the University of Michigan was fired last week after she allegedly said Jewish students were"wealthy and privileged"and did not need support.
Rachel Dawson, who is in charge of the university ' s multicultural initiatives, also purportedly said that the college was"controlled by wealthy Jews"and that"Jewish people have no genetic DNA that would connect them to the land of Israel".
Ms Dawson has denied making those statements or any anti-Semitic comments, according to CNN.

% matched lemmas: 
\end{textsample}
